\chapter{1998:Walter Kohn and John Pople}

\label{chap:chemistry1998}

The Nobel Prize in Chemistry for the year 1998 was awarded jointly to Walter Kohn and John Pople.

\textbf{Motivation:}

for his development of the density-functional theory (one half, Walter Kohn) and for his development of computational methods in quantum chemistry (the other half, John Pople)

\section{Walter Kohn}

\subsection*{Early Life and Education}

Walter Kohn was born in 1923 in Vienna, Austria.

\subsection*{Career and Major Research}

Kohn studied at the University of Toronto (B.A., 1945; M.A., 1946) and earned his Ph.D. from Harvard University in 1948. He held positions at the Carnegie Institute of Technology and the University of California, San Diego, before joining the University of California, Santa Barbara.

\subsection*{Nobel Prize-winning Work}

The work for which Walter Kohn was awarded the Nobel Prize in Chemistry in 1998 was described by the Nobel Committee as: "for his development of the density-functional theory".

He formulated density functional theory, which describes the energy of a system of electrons in terms of the electron density rather than individual wavefunctions, and he devised the Kohn-Sham equations that made the method practical.

\subsection*{Significance and Impact}

Density functional theory became the most widely used computational method for the electronic structure of molecules and solids.

\subsection*{Later Career and Honors}

Kohn continued his research at the University of California, Santa Barbara. He received the 1998 Nobel Prize in Chemistry, and he died in 2016.

\subsection*{Legacy}

The Kohn-Sham method is the foundation of modern computational chemistry and materials science.

\subsection*{Selected Publications}

"Self-Consistent Equations Including Exchange and Correlation Effects" (Physical Review, 1965)

\clearpage

\section{John Pople}

\subsection*{Early Life and Education}

John Pople was born in 1925 in Burnham-on-Sea, England.

\subsection*{Career and Major Research}

Pople studied at the University of Cambridge (B.A., 1946; Ph.D., 1951). He worked at the National Physical Laboratory and Carnegie Mellon University, and from 1986 he was a professor at Northwestern University.

\subsection*{Nobel Prize-winning Work}

The work for which John Pople was awarded the Nobel Prize in Chemistry in 1998 was recognized for his development of computational methods in quantum chemistry.

He created the Gaussian series of computer programs and developed model chemistries that made ab initio quantum chemistry widely accessible.

\subsection*{Significance and Impact}

His computational methods allowed chemists to calculate molecular structures and reaction energies accurately as a matter of routine.

\subsection*{Later Career and Honors}

Pople was a professor at Northwestern University. He shared the 1998 Nobel Prize in Chemistry, and he died in 2004.

\subsection*{Legacy}

The Gaussian software he created remains one of the most used programs in computational chemistry.

\subsection*{Selected Publications}

"Approximate Molecular Orbital Theory" (McGraw-Hill, 1970)

\clearpage

\section*{Chapter Summary}

The 1998 Nobel Prize in Chemistry recognized computational methods in quantum chemistry. Kohn developed density functional theory, while Pople created practical computational methods and software that made such calculations routine.
